% Pursuing Stacks (cotes 134-2 to 134-8): archive cote 134-6, site page 64 (PDF page 65), typescript page 382.
% typescript. Draft compared: ps5.tex:4608-4669.
% AI-transcribed by Claude Opus 5.5 (claude-opus-5-5, adaptive thinking, effort high), Message Batches API, batch ps3 (msgbatch_01XdFmuVVBAcgc7EFxdtCqsj), 2026-11-12; not yet checked by eye.
\scanpage{134-6/64}
\mspage{382}

11.8.

103. It has been over three weeks now I hav'nt been working on the notes.
Most part of the time was spent wandering in the Pyrénées with some friend\hand{s}
(a kind of thing I had'nt been doing since I was a boy), and touring some
other friends living the simple life a\inserted{r}ound there, in the mountains. I was
glad to meet them and \xout{glad} happy to wander and breethe the fresher air of
the mountains - and very happy too after two weeks to be back in the fa-%
miliar surroundings of my home amidst the gentle hills cowered with vine-%
yards ... Yesterday I resumed mathematical work - I had to spend the day
doing scratchwork in order to get back into it, now I feel ready to go on
with the notes. I'll have to finish in the long last with that abelianiza-%
tion story I got into unpremeditatedly - which turns out to be essentially
the same thing as some systematics about (commutative) cohomology and
homology, in the context of ``models'' in (Cat), or in a category $A^\wedge$
(with $A$ in (Cat)). We were out for proving a statement about ``integrato\hand{rs''}
for a\inserted{small} category $B$ being ``abelianizators'' for the dual category \xout{Ax} $A=B^{\hand{$\bullet$}}$ .
The proof I had in mind for this \xout{was} is somewhat indirect via cohomology,
and follows the proof I gave myself a very long time ago (in case $A=\hand{$\Delta$}$ )\hand{,}
that the usual semi-simplicial boundary operations do give the correct
(topos-theoretic) cohomology invariants for any object $X$ in $A^\wedge$ (i.e. any
semisimplicial set), for any locally constant coefficients\unsure{$_{*}$} on $A_{/X}$ .
The idea was to replace $A_{/X}$ by the dual category $(A_{/X})^{o} = {}_{X\hand{$\backslash$}}A^{o}$ (which,
according to a nice result of Quillen, has a homotopy type canonically
isomorphic to the one defined by $A_{/X}$ ), and use the canonical functor
\[ f = \xout{fx}\ (p_X)^{o} : (A_{/X})^{o} \hand{$\longrightarrow$} A^{o} \hand{$\overset{\mathrm{def}}{=} B$}\hand{,} \]
which is a cofibration with discrete fibers, and hence gives rise, for
any \xout{\unsure{sheaf pre}} abelian presheaf $F$ on the category \xout{\illegible{1}}\inserted[typed]{C upstairs}, to an isomor-%
phism
\[ R\Gamma(\over{\illegible{1}}{C}^{\hand{$*$}},F) \simeq R\Gamma(B,f_*(F)) \]
(due to $Rf_*(F)\hand{$\xleftarrow{\ \sim\ }$}f_*(F)$ , as $f_*$ is exact, due to the fact that $f$ is
a cofibration with discrete fibers). We'll get Quillen's result about
the isomorphism $C \simeq C^{o}$ in (Hot), for any object $C$ in (Cat), very
\margin{\uline{VII}} smoothly in part VI, as a result of the asphericity story of part IV.
However, I now realize that the proof of the fact about abelianizators
via Quillen's result and cohomology is rather awkward, as what we're
after\inserted{now} is typically a result on homology, not cohomology - and \deleted{we're}\inserted{I was}
really turning it upside down in order to fit it at all costs into the

\margin{64}
